\documentclass[xelatex,ja=standard,a4paper,10pt]{bxjsarticle}
\geometry{margin=18mm,top=20mm,bottom=22mm,columnsep=7mm}
\usepackage{amsmath,amssymb,mathtools,bm}
\usepackage{graphicx}
\usepackage{booktabs}
\usepackage{multirow}
\usepackage{tabularx}
\usepackage{array}
\usepackage{siunitx}
\usepackage{xcolor}
\usepackage{tikz}
\usetikzlibrary{arrows.meta,positioning,fit,shapes.geometric}
\usepackage{enumitem}
\usepackage{caption}
\usepackage{subcaption}
\usepackage{float}
\usepackage{placeins}
\usepackage{url}
\usepackage[
  colorlinks=true,
  linkcolor=blue!45!black,
  citecolor=blue!45!black,
  urlcolor=blue!55!black
]{hyperref}
\ifpTeX
  \usepackage{pxjahyper}
\fi
\usepackage[nameinlink,noabbrev]{cleveref}

\captionsetup{font=small,labelfont=bf}
\setlength{\emergencystretch}{2em}
\setlength{\parindent}{1em}
\setlength{\parskip}{0pt}
\setlist{nosep,leftmargin=1.8em}
\sisetup{
  detect-all,
  group-separator={,},
  group-minimum-digits=4,
  range-phrase={--},
  range-units=single
}

\newcommand{\pt}{p_{\mathrm{T}}}
\newcommand{\met}{E_{\mathrm{T}}^{\mathrm{miss}}}
\newcommand{\Htt}{H\rightarrow\tau^{+}\tau^{-}}
\newcommand{\Ztt}{Z\rightarrow\tau^{+}\tau^{-}}
\newcommand{\epssig}{\varepsilon_{\mathrm{sig}}}
\newcommand{\epsbkg}{\varepsilon_{\mathrm{bkg}}}
\DeclareSIUnit{\GeV}{GeV}
\DeclareSIUnit{\TeV}{TeV}
\newcommand{\code}[1]{\texttt{\detokenize{#1}}}
\newcommand{\todo}[1]{\textcolor{red!70!black}{[要確認: #1]}}

\crefname{figure}{図}{図}
\crefname{table}{表}{表}
\crefname{section}{節}{節}
\crefname{equation}{式}{式}

\hypersetup{
  pdftitle={再構成されたタウ対による等質量H/Z分離の探索},
  pdfauthor={Ryunosuke Baba}
}


\title{等質量 \(H/Z\rightarrow\tau^{+}\tau^{-}\) 分離レポート\\
\large 論文構成・図表計画}
\author{Ryunosuke Baba}
\date{2026年7月30日}

\begin{document}
\maketitle

\section{論文の中心メッセージ}

truth質量をそれぞれ \(\SI{91.1872}{\GeV}\) と
\(\SI{91.1876}{\GeV}\) に設定した \(H\) と \(Z\) の再構成
\(\tau\)候補を含むMonte Carlo事象を用いた。各データ分割・truth親ボソン
\(\pt\) の20 GeV binごとにH/Z事象数を等しくした後でも、再構成された
\(\tau\)候補、track、PFO、\(\met\) にはH/Z labelを分離する情報が残った。
最終Transformerの固定評価分割ROC AUCは \(0.6088\) であった。

ただし、本解析はスピン相関だけを明示的に抽出した測定ではない。
20 GeV bin内部を含む生成・再構成差を除き切っておらず、再構成
\(\tau\)候補のtruth組成、系統誤差、複数seedも評価していない。
固定評価分割も先行モデルに使用済みである。したがって、結果は
固定partial Monte Carlo標本に対する再構成候補分類の探索的な概念実証として
提示する。

\section{章立て}

\begin{table}[H]
  \centering
  \small
  \begin{tabularx}{\textwidth}{@{}p{0.15\textwidth}X X@{}}
    \toprule
    章 & 書く内容 & 読者に答える問い \\
    \midrule
    要旨
      & sample、\(\pt\) matching、最終AUC、解釈上の限界を数値とともに要約する。
      & 何を行い、何が分かり、どこまで主張できるか。\\
    1. 序論
      & \(\tau\) 崩壊が親粒子のスピン情報を運ぶこと、等質量 \(H/Z\) 比較の動機、
        Transformerを使う理由を説明する。
      & なぜ質量差を消して \(H/Z\) を分けるのか。\\
    2. MCと再構成
      & \(H/Z+j\) 生成、TauDecay、AF3、63.1\% snapshot、CP object処理、
        再構成\(\tau\)候補selection、ntuple検証を書く。
      & 入力事象はどのように作られ、どの事象を使ったか。\\
    3. 解析手法
      & 決定的分割、分割内20 GeV-bin \(\pt\)整合、入力変数、
        前処理、token列、Transformer、HPOの探索・停止・選択規則を定義する。
      & 解析手順をどう固定し、どこまで再現できるか。\\
    4. 結果
      & \(\pt\)整合診断、HPOの完了数と選択結果、40 epoch学習曲線、ROC、
        score、70\%信号効率点、\(\pt\)-bin AUCを示す。
      & 最終分離性能はいくらで、学習は安定しているか。\\
    5. 解釈と制約
      & 粗い親\(\pt\)整合の効果と限界、候補組成、partial sample、
        nominal-only、single seed、評価分割の再利用、MET・trackの未監査点を整理する。
      & 何をスピン効果と呼べず、次に何を検証すべきか。\\
    6. 結論
      & 数値結果と限定付き結論を短く再掲する。
      & この実習から得られた最終的な到達点は何か。\\
    付録
      & 全入力特徴量、model profile、HPO停止条件、software、成果物hashを置く。
      & 本文を切らずに再現性をどう担保するか。\\
    \bottomrule
  \end{tabularx}
  \caption{本文の構成と各章の役割。}
\end{table}

\section{図の計画}

\begin{enumerate}
  \item \textbf{解析フローとtoken表現（新規TikZ図）}: DAODからselection、
    split、\(\pt\)整合、Tensor化、Transformer、validation選択、
    固定評価までを二段で示す。
  \item \textbf{HPOのlearning-rate--dropout散布図}: 21 trialの観測点とtrial 5を示す。
    TPEの点密度を性能面と解釈しない旨を図注へ入れる。
  \item \textbf{40 epoch学習曲線}: epoch 24以降、train lossだけが下がり、
    validation lossとAUCが悪化する過学習を示す。
  \item \textbf{検証/固定評価ROC}: 固定評価AUC \(0.6088\) を主要性能として示す。
    70\%信号効率点は固定評価ROC内部で読み取った記述値であり、
    事前固定閾値ではないと明記する。
  \item \textbf{固定評価score分布}: H/Zの重なりが大きい一方、分布がずれることを示す。
    sigmoid出力の確率calibrationは未評価と明記する。
  \item \textbf{固定評価 \(\pt\)-bin AUC}: class-stratified bootstrap 1000回の68\%区間を示し、
    error barは系統誤差でないと明記する。
\end{enumerate}

本文では個別図を読みやすい大きさで使い、同じ内容を縮小再掲する
summary panelは採用しない。

\section{表の計画}

\begin{enumerate}
  \item H/ZのDAOD数、raw event、selection後event、selection効率。
  \item 再構成\(\tau\)候補とeventのselection条件。
  \item \(\pt\)整合前後およびtrain/validation/testの事象数。
  \item EVENT/TAU/TRACK/PFOの入力と次元、Transformerの最終構成。
  \item HPO範囲とvalidation-onlyで選ばれたtrial 5の設定。
  \item 検証/固定評価のAUC、BCE loss、固定評価ROC上の70\%信号効率点。
  \item 付録にmodel profile、dataset、manifest、checkpointのSHA-256。
\end{enumerate}

\section{論文中で避ける主張}

\begin{itemize}
  \item 「スピン相関だけでAUC \(0.6088\) を得た」とは書かない。
  \item 「20 GeV bin整合で親\(\pt\)分布を完全に同一にした」とは書かない。
  \item 「解析全体でtestを一度しか見ていない」とは書かない。同じtest splitは、
    HPO直後モデルと40 epoch最終モデルのそれぞれで評価済みである。
  \item 再構成\(\tau\)候補selectionの通過率を、真の
    \(\tau_{\mathrm{had}}\tau_{\mathrm{had}}\)信号効率とは呼ばない。
  \item 70\%信号効率点をvalidationで固定したoperating pointとは呼ばない。
  \item bootstrap区間をsystematic uncertaintyと呼ばない。
  \item 63.1\% snapshotを完成sampleまたは実データ性能として扱わない。
\end{itemize}

\end{document}
